% ======================================================================
\section{Electronic models}
\label{sec:models}
% ======================================================================

\subsection{Doped monolayer graphene}
\label{sec:slgmodel}

Near one Dirac point, $h(\bp)=\hbar v_F(\tau p_x\sigma_x+p_y\sigma_y)$ with $\bp$ measured from the Dirac
point, $\hbar v_F=0.658$ eV\,nm ($v_F=10^6$ m/s) and $g=g_sg_v=4$. The bands and spinors are
\begin{equation}
\xi^s_{\bp}=s\hbar v_F|\bp|-E_F,\qquad u^s_{\bp}=\tfrac{1}{\sqrt2}(1,\,s\,e^{\ii\phi_{\bp}}),\qquad s=\pm1,
\end{equation}
with $k_F=\sqrt{4\pi n/g}$ and $E_F=\hbar v_Fk_F$ ($n=10^{12}$ cm$^{-2}$ $\Rightarrow$ $k_F=0.1772$ nm$^{-1}$,
$E_F=116.6$ meV). The overlap that dresses every Coulomb vertex is the chirality factor
\begin{equation}
F^{s's}(\bp,\bp')=|\braket{u^{s'}_{\bp}}{u^{s}_{\bp'}}|^2=\tfrac12\bigl[1+ss'\cos\angle(\bp,\bp')\bigr],
\qquad \textstyle\sum_sF^{s's}=1,
\label{eq:chirality}
\end{equation}
which forbids intraband backscattering and interband forward scattering
(\code{slg\_self\_energy.make\_dirac\_form\_factor}). The sample has only one sheet, so the layer-projected
form factor is $\rho^\ell=\delta_{\ell S}\braket{u}{u'}$: a sample electron has no intermediate state in the tip.

\subsection{Twisted bilayer graphene: continuum model}
\label{sec:tbgmodel}

\paragraph{Hamiltonian.} One valley ($\tau=+$) of the Bistritzer--MacDonald model in the form of
Koshino \emph{et al.}~\cite{BM2011,Koshino2018} as implemented in \texttt{pybandstructure}
(\code{koshino}). In the basis of plane waves $\bk+\bG$ ($\bG$ in the moir\'e reciprocal lattice,
$|\bG|\le G_{\max}|\bG_M|$, $G_{\max}=3$, i.e.\ 37 shells) and of the four orbitals
$(\ell,\sigma)\in\{t,b\}\times\{A,B\}$,
\begin{gather}
H_{\bG\bG'}(\bk)=\delta_{\bG\bG'}
\begin{pmatrix}h_t(\bk+\bG)&0\\0&h_b(\bk+\bG)\end{pmatrix}
+\sum_{j=1}^3\Bigl[\delta_{\bG-\bG',\bm g_j}\begin{pmatrix}0&T_j\\0&0\end{pmatrix}+{\rm h.c.}\Bigr],
\nonumber\\
h_\ell(\bk)=\hbar v\,\bigl[R_\ell(\bk-\bK_\ell)\bigr]\cdot(\tau\sigma_x,-\sigma_y),
\end{gather}
where $R_{t,b}=O(\pm\theta_{\rm TBG}/2)$, $\bK_{1,2}=\frac{|\bG_M|}{\sqrt3}(-\frac{\sqrt3}{2},\pm\frac12)$ are the
moir\'e $\kappa$ points at which the top ($1$) and bottom ($2$) Dirac cones sit, $\bm g_j\in\{0,\bm b_1,\bm b_2\}$
and $T_j=u_{AA}\sigma_0+u_{AB}(\ldots)$ are the three interlayer hoppings with
$|\bG_M|=8\pi\sin(\theta_{\rm TBG}/2)/(\sqrt3a)$, $a=0.246$ nm. Parameters (Ref.~\cite{Wei2025}):
$\theta_{\rm TBG}=1.20^\circ$, $u_{AB}=117$ meV, $u_{AA}=0.7u_{AB}$, $\hbar v=0.657$ eV\,nm, so that
$\alpha=u_{AB}/(\hbar vk_\theta)=0.56$ and the two central bands are 24--25 meV wide and separated from the
remote bands by $\approx44$ meV. Only the $2n_b=8$ bands closest to neutrality are computed
(\code{TBGParams.n\_bands}$=4$). The electrostatic environment differs from Ref.~\cite{Wei2025}, which uses an
isotropic background $\varepsilon_r=10$: here the TBG sits in anisotropic hBN,
$(\varepsilon_\parallel,\varepsilon_z)=(6.41,3)$, with vacuum between its layers (\code{TBGSLGConfig.eps\_par},
\code{eps\_z}, \code{eps\_interlayer}), so plasmon energies are not expected to coincide with that reference.

\paragraph{Model frame vs lab frame.} The model momentum $\bk$ is measured from the moir\'e $\gamma$ point.
The graphene (lab) momentum of the plane wave $\bk+\bG$ in either layer is
\begin{equation}
\bp=\bk+\bG+\bm\Gamma^{\rm lab},\qquad \bK^{\rm lab}_{\rm top}=\bm\Gamma^{\rm lab}+\bK_1,\quad
\bK^{\rm lab}_{\rm bot}=\bm\Gamma^{\rm lab}+\bK_2,
\label{eq:labframe}
\end{equation}
where $\bm\Gamma^{\rm lab}$ (\code{gamma\_point\_tblg}) lies on the $x$ axis and $|\bK^{\rm lab}_{\rm top,bot}|=4\pi/3a$
at angles $\pm\theta_{\rm TBG}/2$ (\code{tbg\_model.model\_frame}; checked in
\code{test\_model\_frame\_dirac\_points}, where the decoupled model has a zero-energy top-layer state at $\bK_1$).
This relation is what the QTM kinematics needs (\cref{sec:qtmtbg}).

\paragraph{Bloch states.} $H(\bk)u_{n\bk}=E_{n\bk}u_{n\bk}$ with components
$u_{n\bk}[4g+2\ell+\sigma]$, $g$ the plane-wave index. The states are stored once
(\code{band\_structure.wavefunctions}) and the \emph{same} gauge is used for the form factors, the self-energy
and the tunneling amplitudes. The chemical potential is fixed by the filling at temperature $T=5$ K
(charge neutrality by default); $\xi_{n\bk}=E_{n\bk}-\mu$.

\paragraph{Layer-resolved form factors.} The density of layer $\ell$ at momentum $\bq+\bG$,
$\rho_\ell(\bq+\bG)=\sum_{\bk}c^\dagger_{\bk}P_\ell e^{-\ii(\bq+\bG)\cdot\br}c_{\bk+\bq}$, has Bloch matrix elements
\begin{equation}
\rho^\ell_{nm}(\bk,\bq;\bG)=\bra{u_{n\bk}}P_\ell\,\hat\tau_{\bG}\ket{u_{m,\bk+\bq}}
\qquad(\code{densityL[l, n, m, k, q, G]}),
\label{eq:formfactor}
\end{equation}
where $P_\ell$ projects on the two sublattices of layer $\ell$ and $\hat\tau_{\bG}$ shifts the plane-wave
index; $\bk+\bq$ is folded back onto the mesh (\code{k\_plus\_q\_ind}) and the corresponding shift of the
plane-wave components is applied by \texttt{pybandstructure}. Completeness:
$\sum_\ell\rho^\ell_{nm}(\bk,0;0)=\delta_{nm}$ (\code{test\_density\_matrix\_completeness}).
Only $\bG=0$ is used for $\Sigma$ (\cref{as:head}).

\paragraph{Doping and Hartree.} For a doped sample the first-shell Hartree potential
$V_H=\sum_{|\bG|=G_M}\frac{2\pi e^2}{\varepsilon_H|\bG|}\delta n(\bG)e^{\ii\bG\cdot\br}$ can be solved
self-consistently (\code{tbg\_model.apply\_hartree}, needs the development version of \texttt{pybandstructure});
the densities $\delta n(\bG)$ are measured from charge neutrality (regularised occupations), and
$G_0$ is then built on the Hartree bands. At neutrality the Hartree potential of the BM model vanishes.

\subsection{Moir\'e Brillouin-zone meshes}

$\bk$ and $\bq$ run over the same $N_s\times N_s$ Wigner--Seitz mesh (\code{n\_sites}), with weights summing to one
and area $A_{\rm BZ}=\frac{\sqrt3}{2}|\bG_M|^2$; the 2D measure is
\begin{equation}
\int\frac{\dd^2q}{(2\pi)^2}f(\bq)\simeq\sum_{\bq}w_{\bq}f(\bq),\qquad w_{\bq}=\frac{A_{\rm BZ}}{(2\pi)^2}\,\text{weight}(\bq)
\qquad(\code{q\_weights\_from\_band\_structure}),
\label{eq:qweights}
\end{equation}
\emph{without} the spin--valley degeneracy (see \cref{sec:degeneracy}); $\bq=0$ is removed ($V$ diverges and its
measure vanishes). The band structure is computed for valley $K$; valley $K'$ follows by time reversal
(\cref{sec:valleys}).
